\documentclass{article}
\usepackage[T1]{fontenc}
\usepackage{lmodern}
\usepackage{graphicx}
\usepackage{amsmath,amssymb}
\usepackage[a4paper,margin=2cm]{geometry}
\usepackage[absolute,overlay]{textpos}
\setlength{\TPHorizModule}{1mm}
\setlength{\TPVertModule}{1mm}
\usepackage[indonesian]{babel}
\usepackage{hyperref}
\setlength{\emergencystretch}{2em}
% unit: Halaman 4

\begin{document}

% === Halaman 4 (python/native) ===

Grup 

• Grup (group) adalah sistem aljabar yang terdiri dari:

 - sebuah himpunan G

 - sebuah operasi biner *

 sedemikian sehingga untuk semua elemen a, b, dan c di dalam G berlaku aksioma 

berikut:

 1.  Closure:  

a * b harus berada di dalam G

 2.  Asosiatif:  

a * (b * c) = (a * b) * c

 3.  Elemen netral: 

terdapat e  G sedemikian sehingga  a * e = e * a = a 

 4.  Elemen invers: 

terdapat a’  G sedemikian sehingga  a * a’ = a’ * a = e 

• Notasi: <G, *>

4

Rinaldi Munir/II4021 Kriptografi

\end{document}
